\section{Historical Stages A--C}\label{app:historical}

Stages A--C precede the prespecified confirmatory experiment reported in Appendix~\ref{app:confirmation}. They are retained because they explain why the final confirmatory analysis was narrowed from a selected-minus-random score to the selected-channel patching score itself.

\subsection{Stage-A prespecified primary comparisons}\label{app:stage-a}

Stage A compared constant template regularization (\texttt{template\_retention\_1}) with template initialization only (\texttt{template\_init}) at 40 epochs. Four $U$ tests were specified before the Stage-A outcomes were inspected and Holm-adjusted as one family. Filter-template similarity and accuracy are reported as descriptive companion outcomes.

\begin{table}[h]
\caption{Complete Stage-A primary family. The positive marginal interval in the first row does not survive Holm adjustment.}
\label{tab:stage-a-primary}
\centering
\small
\setlength{\tabcolsep}{3pt}
\begin{tabular}{llrrrr}
\toprule
Task & Architecture & $\Delta$ template similarity & $\Delta$ acc. (pp) & $\Delta U$ [95\% CI] & Holm $p$ \\
\midrule
single shape & TinyCNN & +0.268 & -0.23 & +0.0345 [0.0018, 0.0672] & 0.1633 \\
single shape & TwoLayerCNN & +0.172 & -0.04 & -0.0092 [-0.0513, 0.0329] & 1.0000 \\
two concepts & TinyCNN & +0.289 & -6.13 & +0.0252 [-0.0306, 0.0811] & 0.9995 \\
two concepts & TwoLayerCNN & +0.198 & -0.39 & -0.0074 [-0.0557, 0.0408] & 1.0000 \\
\bottomrule
\end{tabular}
\end{table}

\subsection{Stage-B epoch-200 condition means}\label{app:stage-b-endpoints}

The table below gives the Stage-B epoch-200 quantities most relevant to the constant-versus-annealed regularization comparison. Means are over ten paired replicates. The full five-condition tables and checkpoint trajectories remain in the analysis artifact.

\begin{table}[h]
\caption{Stage-B epoch-200 means for constant and annealed template regularization.}
\label{tab:stage-b-endpoints}
\centering
\footnotesize
\setlength{\tabcolsep}{3pt}
\begin{tabular}{lllrrr}
\toprule
Task & Architecture & Condition & Accuracy (\%) & Template similarity & $U_{\mathrm{random}}$ \\
\midrule
single shape & TinyCNN & constant $\lambda=1$ & 99.961 & 0.959 & 0.358 \\
single shape & TinyCNN & annealed & 100.000 & 0.690 & 0.335 \\
\addlinespace[2pt]
single shape & TwoLayerCNN & constant $\lambda=1$ & 99.922 & 1.000 & 0.067 \\
single shape & TwoLayerCNN & annealed & 99.961 & 0.942 & 0.060 \\
\addlinespace[2pt]
two concepts & TinyCNN & constant $\lambda=1$ & 98.008 & 0.947 & 0.657 \\
two concepts & TinyCNN & annealed & 99.219 & 0.580 & 0.718 \\
\addlinespace[2pt]
two concepts & TwoLayerCNN & constant $\lambda=1$ & 99.141 & 0.999 & 0.061 \\
two concepts & TwoLayerCNN & annealed & 99.531 & 0.907 & 0.003 \\
\bottomrule
\end{tabular}
\end{table}

On \texttt{two\_concepts}/TinyCNN, patched-run accuracy on the counterfactual target is 93.61\% under constant regularization and 90.64\% under annealed regularization even though $U_{\mathrm{random}}$ is higher under annealing. This discrepancy motivated separate reporting of the selected-channel patching score and selected-minus-control scores.

\subsection{Stage-C decomposition of the original selected-minus-random score}\label{app:stage-c-decomposition}

Stage C reused the 80 epoch-200 Stage-B checkpoints from the annealed- and constant-regularization conditions and evaluated three validation-only channel-selection rules at $k\in\{1,2,4,8\}$. The central TinyCNN result is reproduced below. Each entry is annealed minus constant regularization; intervals are marginal paired 95\% $t$ intervals over ten replicates. Stage C is post hoc and supplies no confirmatory test.

{
\small
\setlength{\tabcolsep}{3pt}
\begin{longtable}{llrrrr}
\caption{Stage-C \texttt{two\_concepts}/TinyCNN decomposition.}\label{tab:stage-c-decomp}\\
\toprule
Channel selection & $k$ & $\Delta F_{sel}$ & $\Delta F_{rand}$ & $\Delta U$ & $\Delta$ CF acc. \\
\midrule
\endfirsthead
\toprule
Channel selection & $k$ & $\Delta F_{sel}$ & $\Delta F_{rand}$ & $\Delta U$ & $\Delta$ CF acc. \\
\midrule
\endhead
\bottomrule
\endfoot
activation difference & 1 & -0.375 & -0.040 & -0.335 & -0.255 \\
activation difference & 2 & -0.163 & -0.062 & -0.101 & -0.185 \\
activation difference & 4 & -0.022 & -0.083 & +0.061 & -0.030 \\
activation difference & 8 & +0.006 & -0.033 & +0.040 & +0.013 \\
AUROC & 1 & -0.332 & -0.040 & -0.292 & -0.230 \\
AUROC & 2 & -0.188 & -0.062 & -0.125 & -0.218 \\
AUROC & 4 & -0.022 & -0.083 & +0.061 & -0.031 \\
AUROC & 8 & +0.007 & -0.033 & +0.040 & +0.015 \\
validation patching & 1 & -0.277 & -0.040 & -0.237 & -0.194 \\
validation patching & 2 & -0.122 & -0.062 & -0.060 & -0.147 \\
validation patching & 4 & -0.015 & -0.083 & +0.067 & -0.021 \\
validation patching & 8 & +0.005 & -0.033 & +0.038 & +0.013 \\
\end{longtable}
\normalsize
}

When channels are ranked by activation difference, $\Delta F_{sel}(4)=-0.022$ with 95\% interval $[-0.054,0.010]$, while $\Delta F_{rand}(4)=-0.083$ with interval $[-0.097,-0.068]$. The positive $\Delta U(4)=+0.061$ is therefore partly explained by the change in the random-channel baseline. At $k=8$, the selected-channel patching score itself becomes slightly positive: $+0.006$ $[0.004,0.009]$. AUROC and validation-patching channel selection show the same broad pattern.

For \texttt{two\_concepts}/TwoLayerCNN, the selected-channel patching score at $k=4$ is negative under all three channel-selection rules (approximately -0.125, -0.121, and -0.087 for activation difference, AUROC, and validation patching, respectively). Under validation-patching selection the score becomes positive at $k=8$ ($+0.0215$ with marginal interval $[0.0049,0.0381]$). These outcomes motivated a confirmatory hypothesis based on changes across the number of patched channels rather than on a fixed sign at one value of $k$.

\subsection{Stage-C validation-energy-matched control}\label{app:stage-c-energy}

A later post hoc control chose alternative channel sets of the same size whose validation replacement energy was closest to that of the selected set. All 80 checkpoint evaluations completed. Matching quality depends strongly on channel count: single-channel sets can have large relative mismatch, while $k=4$ and especially $k=8$ are generally tightly matched in the key TinyCNN setting. The selected-channel annealed-minus-constant contrast is unchanged by baseline construction and remains the direct quantity used to motivate Stage D. The energy-matched selected-minus-control score is therefore reported as a secondary analysis.

\subsection{Transition to Stage D}\label{app:transition}

Stages A--C established that constant template regularization preserves high filter-template similarity and that prediction can improve after the regularizer is annealed to zero. They also showed that the annealed-minus-constant patching comparison changes with the number of patched channels. Because the original $U$ sign change combined selected-channel and baseline patching scores, the Stage-D protocol was specified around the selected-channel score before new replicates were generated. Appendix~\ref{app:confirmation} reports that confirmation and the subsequent 1,200-model secondary sensitivity study.